\section*{Hoofdstuk \ref{ch:b1:angular-momentum} --- Impulsmoment}

\begin{solution}{exo:b1:angular-momentum:1}
$Fd = 50 \times 0.30 = \qty{15}{N.m}$; $Fd\sin 60^\circ = \qty{13}{N.m}$;
langs de steel gaat de werklijn door de moer: $0$.
\end{solution}

\begin{solution}{exo:b1:angular-momentum:2}
$L = mrv = 5.97\times10^{24} \times 1.50\times10^{11} \times 2.98\times10^4 =
\qty{2.7e40}{kg.m^2/s}$; $\dd A/\dd t = L/2m = rv/2 = \qty{2.2e15}{m^2/s}$.
\end{solution}

\begin{solution}{exo:b1:angular-momentum:3}
$r = \ell\sin 30^\circ = \qty{0.50}{m}$, $\omega = \qty{3.4}{rad/s}$:
$L_z/m = r^2\omega = \qty{0.84}{m^2/s}$. De werklijn van de spankracht
gaat door het draaipunt (\omterm{def:b1:angular-momentum:moment}{moment} nul om elke as erdoorheen); het
gewicht is evenwijdig aan de verticale as (\omterm{def:b1:angular-momentum:moment}{moment} nul daarom): $L_z$
blijft behouden.
\end{solution}

\begin{solution}{exo:b1:angular-momentum:4}
$J_1 = 4.0 + 2 \times 3.0 \times 0.64 = \qty{7.8}{kg.m^2}$, $J_2 = 4.0 +
2 \times 3.0 \times 0.04 = \qty{4.2}{kg.m^2}$; $\omega_2 = \omega_1J_1/J_2 =
1.85\omega_1$: \qty{3.7}{} omwentelingen per seconde. Met $E_k = L^2/2J$ is de verhouding
$J_1/J_2 = 1.85$; de spieren, die de \omterm{def:b1:angular-momentum:moment}{armen} naar binnen trekken tegen de
middelpuntvliedende neiging in, verrichten de arbeid.
\end{solution}

\begin{solution}{exo:b1:angular-momentum:5}
$L_\Delta = m\ell^2\dot\theta$; spankracht: \omterm{def:b1:angular-momentum:moment}{moment} nul (snijdt de as);
gewicht: $-mg\ell\sin\theta$ (\omterm{def:b1:angular-momentum:moment}{arm} $\ell\sin\theta$, terugdrijvend).
$m\ell^2\ddot\theta = -mg\ell\sin\theta$.
\end{solution}

\begin{solution}{exo:b1:angular-momentum:6}
$r_pv_p = r_av_a$: $v_p/v_a = r_a/r_p = 1.034$; met $(v_p + v_a)/2 \approx
\qty{29.8}{km/s}$: $v_p = \qty{30.3}{km/s}$, $v_a = \qty{29.3}{km/s}$.
\end{solution}

\begin{solution}{exo:b1:angular-momentum:7}
$L = mrv$ blijft behouden: $v_2 = v_1r_1/r_2 = \qty{4.0}{m/s}$, $\omega_2 = v_2/r_2
= \qty{16}{rad/s}$ (vanaf \qty{4}{rad/s}). $E_1 = \qty{0.40}{J}$, $E_2 =
\qty{1.6}{J}$: het trekken levert \qty{1.2}{J}. De spankracht is radiaal, maar de
verplaatsing van de puck heeft nu een radiale component (zij spiraalt naar binnen): de arbeid
$\vect T\cdot\dd\vect\ell$ is niet langer nul.
\end{solution}

\begin{solution}{exo:b1:angular-momentum:8}
\omterm{def:b1:angular-momentum:moment}{Momenten} om het draaipunt (lengtes vanaf het draaipunt): kind $30 \times 1.5 =
45$ (links); het gewicht van de plank in haar midden, $0.5$ naar rechts: $20 \times
0.5 = 10$ (rechts); tweede kind op $x$ naar rechts: $45x$. Evenwicht:
$45 = 10 + 45x$, $x = \qty{0.78}{m}$. Het draaipunt draagt het totale gewicht,
$95 \times 9.81 = \qty{932}{N}$.
\end{solution}

\begin{solution}{exo:b1:angular-momentum:9}
Massa's: $m_1a = m_1g - T_1$, $m_2a = T_2 - m_2g$; katrol: $J\dot\omega =
(T_1 - T_2)R$ met $a = R\dot\omega$. Optellen: $a = (m_1 - m_2)g/(m_1 + m_2 +
J/R^2) = 9.81/(3.0 + 1.0) = \qty{2.5}{m/s^2}$ (tegen \qty{3.3}{m/s^2});
$T_1 = m_1(g - a) = \qty{14.7}{N}$, $T_2 = m_2(g + a) = \qty{12.3}{N}$ ---
ongelijk, want de katrol heeft een netto \omterm{def:b1:angular-momentum:moment}{moment} nodig.
\end{solution}

\begin{solution}{exo:b1:angular-momentum:10}
$r_1v_1 = r_2v_{\perp 2}$: $v_{\perp 2} = 40 \times 0.5/10 = \qty{2.0}{km/s}$.
De \omterm{cor:b1:angular-momentum:central}{perkenwet} geeft alleen de loodrechte component; voor de radiale is
het behoud van energie nodig (\cref{ch:b1:central-forces}).
\end{solution}

\begin{solution}{exo:b1:angular-momentum:11}
$J\omega$ constant met $J \propto R^2$: $T' = T(R'/R)^2 = 25 \times 86400
\times (10^4/7\times10^8)^2 = \qty{4.4e-4}{s}$. Snelheid aan de evenaar $2\pi R'/T'
= \qty{1.4e8}{m/s}$, de helft van de lichtsnelheid: de ineenstorting kan niet
al het \omterm{def:b1:angular-momentum:L}{impulsmoment} behouden --- veel gaat verloren (en de relativiteit komt erbij).
\end{solution}

\begin{solution}{exo:b1:angular-momentum:12}
$L_z = mr^2\omega$, en $r$ verandert: $\dd L_z/\dd t = 2mr\dot r\omega \neq 0$.
Dat moet het \omterm{def:b1:angular-momentum:moment}{moment} van de dwarse reactie $N$ van de staaf zijn: $rN =
2mr\dot r\omega$, dus $N = 2m\dot r\omega$ --- de staaf duwt de kraal zijwaarts
terwijl zij naar buiten glijdt (de coriolisterm van \cref{ch:b1:non-inertial-frames}).
\end{solution}

\begin{solution}{pb:b1:angular-momentum:1}
\textbf{1.} $m\omega^2r = GMm/r^2$: $\omega = \sqrt{GM/r^3} = \sqrt{3.98\times
10^{14}/5.66\times10^{25}} = \qty{2.65e-6}{rad/s}$, periode $2\pi/\omega =
\qty{2.37e6}{s} = \qty{27.4}{d}$.

\textbf{2.} $L_{\text{baan}} = 7.35\times10^{22} \times 1.475\times10^{17}
\times 2.65\times10^{-6} = \qty{2.9e34}{kg.m^2/s}$.

\textbf{3.} $v = \sqrt{GM/r}$, $L = mrv = m\sqrt{GMr}$; $\dd L/\dd r =
\tfrac12 m\sqrt{GM/r} = L/2r$.

\textbf{4.} $J = 0.33 \times 5.97\times10^{24} \times 4.06\times10^{13} =
\qty{8.0e37}{kg.m^2}$; $\Omega = \qty{7.29e-5}{rad/s}$; $L_{\text{spin}} =
\qty{5.8e33}{kg.m^2/s}$, een vijfde van $L_{\text{baan}}$.

\textbf{5.} $J_m \approx 0.4 \times 7.35\times10^{22} \times 3.0\times10^{12} =
\qty{8.9e34}{kg.m^2}$, maal $\omega$: \qty{2.4e29}{kg.m^2/s}, $10^{-5}$ van
dat van de \omterm{def:b1:point-kinematics:frame}{baan}: verwaarloosbaar.

\textbf{6.} De aantrekking van de zon grijpt aan in het massamiddelpunt van het stelsel
(in eerste orde geen \omterm{def:b1:angular-momentum:moment}{moment} daarom); haar getijdeneffect op het paar is
klein: het totale $L$ blijft behouden met de hier vereiste nauwkeurigheid.

\textbf{7.} De nabije vloedberg, die vóór de lijn aarde--maan uitloopt, wordt
door de maan aangetrokken met een kracht waarvan de tangentiële component de
draaiing van de aarde tegenwerkt.

\textbf{8.} Een negatief \omterm{def:b1:angular-momentum:moment}{moment} op de draaiing: de aarde vertraagt, de dag
wordt langer.

\textbf{9.} De vloedberg trekt de maan vooruit: een positief \omterm{def:b1:angular-momentum:moment}{moment} op de
\omterm{def:b1:point-kinematics:frame}{baan}, $L_{\text{baan}}$ groeit.

\textbf{10.} $L_{\text{baan}} \propto \sqrt r$ groeit, dus groeit $r$; $v \propto
1/\sqrt r$ neemt af.

\textbf{11.} Het vooruit trekken verricht positieve arbeid en verhoogt de energie
van de maan; een hogere \omterm{def:b1:point-kinematics:frame}{baan} is een tragere \omterm{def:b1:point-kinematics:frame}{baan}: de energie gaat naar hoogte
(potentieel), meer dan de \omterm{thm:b1:work-and-energy:ke}{kinetische energie} die verloren gaat.

\textbf{12.} Getijdenkoppeling van de aarde: dag gelijk aan maand, geen voorlopende
vloedberg, geen \omterm{def:b1:angular-momentum:moment}{koppel}.

\textbf{13.} $\Delta\Omega/\Omega = -\qty{2.3e-3}{s}/\qty{86164}{s} =
\num{-2.7e-8}$ per eeuw.

\textbf{14.} $\Delta L_{\text{spin}} = L_{\text{spin}}\Delta\Omega/\Omega =
\qty{-1.55e26}{kg.m^2/s}$ per eeuw.

\textbf{15.} $\Delta L_{\text{baan}} = +\qty{1.55e26}{kg.m^2/s}$ per eeuw.

\textbf{16.} $\Delta r = 2r\,\Delta L/L = 2 \times 3.84\times10^8 \times 1.55
\times10^{26}/2.9\times10^{34} = \qty{4.1}{m}$ per eeuw, \qty{4.1}{cm/yr}:
dicht bij de gemeten \qty{3.8}{cm/yr} (het langer worden van de dag heeft ook
een niet-getijdengebonden deel).

\textbf{17.} $\Delta T/T = \tfrac32\Delta r/r = 1.6\times10^{-8}$: $2.37\times
10^6 \times 1.6\times10^{-8} = \qty{0.04}{s}$ per eeuw.

\textbf{18.} $\Delta E_{\text{rot}} = L_{\text{spin}}\Delta\Omega = 5.8\times10^{33}
\times 7.29\times10^{-5} \times (-2.7\times10^{-8}) = \qty{-1.1e22}{J}$ per
eeuw.

\textbf{19.} $\Delta E_{\text{baan}} = GMm\Delta r/2r^2 = 2.93\times10^{37}
\times 4.1/(2 \times 1.475\times10^{17}) = \qty{+4.1e20}{J}$ per eeuw.

\textbf{20.} Netto $-1.1\times10^{22} + 4\times10^{20} \approx \qty{-1.06e22}{J}$
per eeuw: $3.4\times10^{12}\ \unit{W}$, als warmte gedissipeerd door getijdenwrijving
in de oceanen en in de vaste aarde.

\textbf{21.} Een vijfde van het vermogen van de mensheid, een veertiende van de
geothermische stroom --- een kleine maar blijvende kachel.

\textbf{22.} $L_{\text{tot}} = 3.5\times10^{34}$. Bij $r = 5.5\times10^8$:
$\omega = \sqrt{GM/r^3} = \qty{1.55e-6}{rad/s}$, $J\omega = 1.2\times10^{32}$,
$m\sqrt{GMr} = 7.35\times10^{22} \times 4.68\times10^{11} = 3.44\times10^{34}$;
som $3.45\times10^{34}$: consistent. Periode $2\pi/\omega = \qty{4.1e6}{s}
= \qty{47}{d}$.

\textbf{23.} De getijden van de zon blijven de aarde afremmen tot onder het
baantempo van de maan, zodat het \omterm{def:b1:angular-momentum:moment}{koppel} aarde--maan dan zou omkeren; bovendien
zal de zon allang zijn geëvolueerd (de tijdschaal bedraagt tientallen
miljarden jaren).

\textbf{24.} De aarde wekte getijden op de maan op, die haar draaiing afremden
tot één zijde naar ons gekeerd bleef; omdat zij lichter is en dicht bij een zwaardere
partner staat, ging de koppeling van de maan veel sneller.

\textbf{25.} Behouden: het totale \omterm{def:b1:angular-momentum:L}{impulsmoment}. Uitgewisseld: \omterm{def:b1:angular-momentum:L}{impulsmoment}
van de draaiing van de aarde naar de \omterm{def:b1:point-kinematics:frame}{baan} van de maan, ongeveer \qty{1.6e26}{SI}
per eeuw. Gedissipeerd: \omterm{thm:b1:work-and-energy:em}{mechanische energie}, enkele terawatt, als
getijdenwarmte.
\end{solution}
